\documentclass[conference]{IEEEtran}

\usepackage{float}
\usepackage{cite}
\usepackage{caption}
\usepackage{subcaption}
\usepackage{amsmath,amssymb,amsfonts,amstext}
% \usepackage{algorithmic}
\usepackage[dvipdfmx]{graphicx}
\usepackage{textcomp}
\usepackage{xcolor}
\usepackage{colortbl}
\usepackage[T1]{fontenc}
\usepackage{mdframed}
\usepackage{algorithm, algpseudocode} % texlive-science
\usepackage[braket, qm]{qcircuit}
\usepackage{newtxtext}
\usepackage[normalem]{ulem}
\usepackage{nameref}
\usepackage{tikz}
\usetikzlibrary{arrows.meta, shapes, positioning, calc, fit}

\input{img/defs}

\makeatletter
\newcommand{\sublabel}[1]{\protected@edef\@currentlabel{\thefigure(\thesubfigure)}\label{#1}}
\makeatother
\captionsetup{subrefformat=parens}
\algrenewcommand\algorithmicindent{0.5em}
\newtheorem{definition}{Definition}
\newtheorem{example}{Example}
\newcommand{\thickhline}{%
    \noalign {\ifnum 0=`}\fi \hrule height 1pt
    \futurelet \reserved@a \@xhline
}
%\def\BibTeX{{\rm B\kern-.05em{\sc i\kern-.025em b}\kern-.08em
%    T\kern-.1667em\lower.7ex\hbox{E}\kern-.125emX}}
% You may prefer to use the caption package instead of the deprecated \@captype hack below.
\newcommand{\figcaption}[1]{\def\@captype{figure}\caption{#1}}
\newcommand{\tblcaption}[1]{\def\@captype{table}\caption{#1}}

\newcommand{\Del}[1]{\textcolor{red}{\sout{#1}}}
\newcommand{\Add}[1]{\textcolor{blue}{\emph{#1}}}

\newcount\NGCcolor
\NGCcolor=1
\newcount\NGCcolortwo
\NGCcolortwo=1
%\NGCcolor=0
\ifnum\NGCcolortwo<1
  \newcommand{\blueoutt}[1]{\textcolor{blue}{#1}}  
\else
  \newcommand{\blueoutt}[1]{#1}
\fi

\ifnum\NGCcolor<1
  \newcommand{\redout}[1]{\textcolor{red}{#1}}
  \newcommand{\blueout}[1]{\textcolor{blue}{#1}}
  \newcommand{\greenout}[1]{\textcolor{green}{#1}}
  \newcommand{\cyannout}[1]{\textcolor{cyan}{#1}}
  \newcommand{\magentaout}[1]{\textcolor{magenta}{#1}}
  \newcommand{\yellowout}[1]{\textcolor{yellow}{\bf #1}}
  \newcommand{\algblueout}[1]{\leavevmode\color{blue}{#1}}
\else
  \newcommand{\redout}[1]{#1}
  \newcommand{\blueout}[1]{#1}
  \newcommand{\greenout}[1]{#1}
  \newcommand{\cyannout}[1]{#1}
  \newcommand{\magentaout}[1]{#1}
  \newcommand{\yellowout}[1]{#1}
  \newcommand{\algblueout}[1]{#1}
\fi

% End preamble
